% Part: normal-modal-logic
% Chapter: filtrations
% Section: examples-of-filtrations

\documentclass[../../../include/open-logic-section]{subfiles}

\begin{document}

\olfileid[fr]{nml}{fil}{exf}

\olsection{Exemples de filtrations}

Nous n'avons pas encore montré qu'il existe des filtrations. Pourtant,
pour tout modèle $\mModel{M}$, il existe au moins une filtration par
$\Gamma$. Nous définirons deux filtrations extrêmes : la plus fine et
la plus grossière, qui peuvent coïncider. Les filtrations d'un même modèle
diffèrent par leur relation d'accessibilité (puisque
la \olref[fil]{defn:filtration} fixe directement $W^*$ et~$V^*$). Dans la
filtration la plus fine, aussi peu de mondes que possible sont reliés ;
dans la plus grossière, autant que possible le sont.

\begin{defn}
  Si $\Gamma$ est clos par sous-formules, la filtration \emph{la plus
  fine} $\mModel{M^*}$ d'un modèle $\mModel{M}$ est définie par :
  \[
  R^*[u][v] \quad \text{si et seulement si} \quad \exists u'\in [u] \;
  \exists v' \in [v] : Ru'v'.
  \]
\end{defn}

\begin{prop}\ollabel{prop:finest}
  La filtration la plus fine $\mModel{M^*}$ est bien une filtration.
\end{prop}

\begin{proof}
  Vérifions que la relation $R^*$ ainsi définie satisfait
  la \olref[fil]{defn:filtration}\olref[fil]{defn:filtration-R}, en
  examinant successivement les trois conditions.

  Si $Ruv$, comme $u \in [u]$ et $v \in [v]$, on a aussi
  $R^*[u][v]$ ; \olref[fil]{defn:filtration-R1} est donc satisfaite.

  \iftag{prvBox}{\iftag{probBox}{Nous laissons en exercice la
      vérification de \olref[fil]{defn:filtration-R2}.}{Pour
        \olref[fil]{defn:filtration-R2}, supposons
        $\Box!A \in \Gamma$, $R^*[u][v]$ et
        $\mSat{M}{\Box!A}[u]$. Par définition de $R^*$, il existe
        $u' \equiv u$ et $v' \equiv v$ tels que $Ru'v'$.
        Comme $u$ et $u'$ s'accordent sur $\Gamma$, on a aussi
        $\mSat{M}{\Box!A}[u']$, d'où $\mSat{M}{!A}[v']$. Puisque
        $\Gamma$ est clos par sous-!!{formula}s, $v$ et $v'$ s'accordent
        sur $!A$ ; donc $\mSat{M}{!A}[v]$, comme voulu.}}{}

  \iftag{prvDiamond}{\iftag{probDiamond}{Nous laissons en exercice la
      vérification de \olref[fil]{defn:filtration-R3}.}{Pour vérifier
        \olref[fil]{defn:filtration-R3}, supposons $\Diamond!A \in
        \Gamma$, $R^*[u][v]$ et $\mSat{M}{!A}[v]$. Par définition de
        $R^*$, il existe $u' \equiv u$ et $v' \equiv v$ tels que
        $Ru'v'$. Comme $v$ et $v'$ s'accordent sur~$\Gamma$ et que
        $\Gamma$ est clos par sous-!!{formula}s, on a aussi
        $\mSat{M}{!A}[v']$, d'où
        $\mSat{M}{\Diamond !A}[u']$. Comme $u$ et $u'$ s'accordent
        également sur $\Gamma$, on a
        $\mSat{M}{\Diamond !A}[u]$.}}{}
\end{proof}

\begin{probtag}{probBox,probDiamond}
  Complétez la démonstration de la \olref[nml][fil][exf]{prop:finest}.
\end{probtag}

\begin{defn}
  Si $\Gamma$ est clos par sous-formules, la filtration \emph{la plus
  grossière}~$\mModel{M^*}$ d'un modèle~$\mModel{M}$ est définie en
  posant $R^*[u][v]$ si et seulement si
  \iftag{notprvBox,notprvDiamond}{la condition suivante est
    satisfaite :}{\emph{les deux} conditions suivantes sont satisfaites :}
  \begin{tagenumerate}{prvBox,prvDiamond}
  \tagitem{prvBox}{\ollabel{defn:coarsest-Box}Si $\Box!A \in \Gamma$
    et $\mSat{M}{\Box!A}[u]$, alors
    $\mSat{M}{!A}[v]$\iftag{prvDiamond}{ ;}{.}}{}
  \tagitem{prvDiamond}{\ollabel{defn:coarsest-Diamond}Si $\Diamond!A
    \in \Gamma$ et $\mSat{M}{!A}[v]$, alors
    $\mSat{M}{\Diamond!A}[u]$.}{}
  \end{tagenumerate}
\end{defn}

\begin{prop}
  La filtration la plus grossière $\mModel{M^*}$ est bien une
  filtration.
\end{prop}

\begin{proof}
  D'après la définition de $R^*$, il ne reste à vérifier que
  l'implication de $Ruv$ vers $R^*[u][v]$. Supposons donc
  $Ruv$. \iftag{prvBox}{Supposons $\Box!A \in \Gamma$ et
    $\mSat{M}{\Box!A}[u]$ ; on a alors évidemment
    $\mSat{M}{!A}[v]$\iftag{prvDiamond}{, ce qui satisfait
      \olref{defn:coarsest-Box}. }{.}}{}\iftag{prvDiamond}{Supposons
    $\Diamond!A \in \Gamma$ et $\mSat{M}{!A}[v]$. Alors
    $\mSat{M}{\Diamond!A}[u]$ puisque~$Ruv$\iftag{prvBox}{, ce qui
      satisfait \olref{defn:coarsest-Diamond}}{}.}{}
\end{proof}

\begin{ex}
  Soit $W = \PosInt$, soit $Rnm$ si et seulement si $m = n + 1$, et
  soit $V(p) = \Setabs{2n}{n \in
    \PosInt}$. Le modèle $\mModel{M} = \tuple{W, R, V}$ est représenté
  dans la \olref{fig:ex-filtration}. Ses mondes sont $1$, $2$, etc. ;
  chaque monde accède à un seul autre monde, son successeur, et $p$
  est vraie exactement aux nombres pairs.

  \begin{figure}
    \centering
    \begin{tikzpicture}[modal]
      \node[world] (1) [label=below:\mFalse{p}] {$1$};
      \node[world] (2) [label=below:\mTrue{p}, right=of 1] {$2$};
      \node[world] (3) [label=below:\mFalse{p}, right=of 2] {$3$};
      \node[world] (4) [label=below:\mTrue{p}, right=of 3] {$4$};
      \node[phantom] (5) [right=of 4] {};
      \draw[->] (1) to (2);
      \draw[->] (2) to (3);
      \draw[->] (3) to (4);
      \draw[dotted] (4) to (5);
    \end{tikzpicture}

    \begin{tikzpicture}[modal]
      \node[world] (1) [label=below:\mFalse{p}] {$[1]$};
      \node[world] (2) [label=below:\mTrue{p}, right=of 1] {$[2]$};
      \draw[->,bend left] (1) to (2);
      \draw[->,bend left] (2) to (1);

      \node[world] (3) [label=below:\mFalse{p},right=of 2] {$[1]$};
      \node[world] (4) [label=below:\mTrue{p}, right=of 3] {$[2]$};
      \draw[->,bend left] (3) to (4);
      \draw[->,bend left] (4) to (3);
      \draw[->,reflexive above] (4) to (4);
    \end{tikzpicture}
    \caption{Un modèle infini et ses filtrations.}
    \ollabel{fig:ex-filtration}
  \end{figure}

  Soit maintenant $\Gamma$ l'ensemble des sous-!!{formula}s de~$\Box p \lif p$,
  à savoir $\{p, \Box p, \Box p \lif p\}$. $p$ est vraie exactement
  aux nombres pairs, tandis que $\Box p$ est vraie exactement aux
  nombres impairs. Ainsi, $\Box p \lif p$ est vraie exactement aux
  nombres pairs. Autrement dit, tout nombre impair rend $\Box p$
  vraie, mais $p$ et $\Box p \lif p$ fausses ; tout nombre pair rend
  $p$ et $\Box p \lif p$ vraies, mais $\Box p$ fausse. On a donc
  $W^* = \{ [1], [2] \}$, où $[1] = \{1, 3, 5, \dots\}$ et
  $[2] = \{2, 4, 6, \dots\}$. Puisque $2 \in V(p)$,
  $[2] \in V^*(p)$ ; puisque $1 \notin V(p)$,
  $[1] \notin V^*(p)$. Par conséquent, $V^*(p) = \{[2]\}$.

  Sur $W^*$, la relation d'accessibilité doit contenir
  $\tuple{[1], [2]}, \tuple{[2],[1]}$ : puisque $R12$,
  on doit avoir $R^*[1][2]$ par la
  \olref[fil]{defn:filtration}\olref[fil]{defn:filtration-R1} ; et,
  puisque $R23$, on doit avoir $R^*[2][3]$, avec $[3]=[1]$. Elle ne
  peut pas contenir $\tuple{[1],[1]}$ : sinon, on aurait
  $R^*[1][1]$, $\mSat{M}{\Box p}[1]$ mais
  $\mSat/{M}{p}[1]$, en contradiction avec
  \olref[fil]{defn:filtration-R2}. Rien n'impose ni n'interdit
  $R^*[2][2]$. Il y a donc deux filtrations possibles de~$\mModel{M}$,
  correspondant aux deux relations d'accessibilité
  \[
  \{\tuple{[1],[2]}, \tuple{[2],[1]}\} \text{ et }
  \{\tuple{[1],[2]}, \tuple{[2],[1]}, \tuple{[2],[2]}\}.
  \]
  Dans les deux cas, $p$ et $\Box p \lif p$ sont fausses et
  $\Box p$ est vraie en~$[1]$ ; $p$ et $\Box p \lif p$ sont vraies et
  $\Box p$ est fausse en~$[2]$.
\end{ex}

\begin{prob}
  Considérez le modèle suivant $\mModel{M} = \tuple{W, R, V}$, où
  $W = \Setabs{0\sigma}{\sigma \in \Bin^*}$ est l'ensemble des suites
  de $0$ et de~$1$ commençant par~$0$, où $R\sigma\sigma'$ si et
  seulement si $\sigma' = \sigma 0$ ou $\sigma' = \sigma 1$, et où
  $V(p) = W \cap \Setabs{\sigma
    0}{\sigma \in \Bin^*}$ et
  $V(q) = W \cap \Setabs{\sigma 1}{\sigma \in
    \Bin^* \setminus \{1\}}$. Voici une représentation :
  \begin{center}
    \begin{tikzpicture}[modal,
        node distance=2cm
      ]

      \node[world] (0) [label={below:\mTrue{p}\\\mFalse{q}},font=\tiny] {$0$};
      \node[world] (00) [label={below:\mTrue{p}\\\mFalse{q}},
        above right=of 0,font=\tiny] {$00$};
      \node[world] (000) [label={below:\mTrue{p}\\\mFalse{q}},
        right=of 00, yshift=.8cm, font=\tiny] {$000$};
      \node[phantom] (0000) [right=of 000, yshift=.5cm] {};
      \node[phantom] (0001) [right=of 000, yshift=-.5cm] {};
      \node[world] (001) [label={below:\mFalse{p}\\\mTrue{q}},
        right=of 00, yshift=-.8cm, font=\tiny] {$001$};
      \node[phantom](0010) [right=of 001, yshift=.5cm] {};
      \node[phantom](0011) [right=of 001, yshift=-.5cm] {};
      \node[world] (01) [label={below:\mFalse{p}\\\mTrue{q}},
        below right=of 0,font=\tiny] {$01$};
      \node[world] (010) [label={below:\mTrue{p}\\\mFalse{q}},
        right=of 01, yshift=.8cm,font=\tiny] {$010$};
      \node[phantom] (0100) [right=of 010, yshift=.5cm] {};
      \node[phantom] (0101) [right=of 010, yshift=-.5cm] {};
      \node[world] (011) [label={below:\mFalse{p}\\\mTrue{q}},
        right=of 01, yshift=-.8cm,font=\tiny] {$011$};
      \node[phantom] (0110) [right=of 011, yshift=.5cm] {};
      \node[phantom] (0111) [right=of 011, yshift=-.5cm] {};

      \draw[->] (0) to (00);
      \draw[->] (0) to (01);
      \draw[->] (00) to (000);
      \draw[dotted] (000) to (0000);
      \draw[dotted] (000) to (0001);
      \draw[->] (00) to (001);
      \draw[dotted] (001) to (0010);
      \draw[dotted] (001) to (0011);
      \draw[->] (01) to (010);
      \draw[dotted] (010) to (0100);
      \draw[dotted] (010) to (0101);
      \draw[->] (01) to (011);
      \draw[dotted] (011) to (0110);
      \draw[dotted] (011) to (0111);
    \end{tikzpicture}
  \end{center}
  On a $\mSat/{M}{\Box(p \lor q) \lif (\Box p \lor \Box q)}[w]$
  pour tout~$w$.

  Soit $\Gamma$ l'ensemble des sous-!!{formula}s de~$\Box(p \lor q) \lif
  (\Box p \lor \Box q)$. Quels sont $W^*$ et~$V^*$ ? Quelle est la
  relation d'accessibilité de la filtration la plus fine de~$\mModel{M}$ ?
  Et de la plus grossière ?
\end{prob}
\begin{explain}
La filtration la plus fine conserve les flèches héritées
du modèle. La plus grossière ajoute toutes les flèches que
permettent les conditions de filtration. Elles coïncident
lorsque les flèches héritées épuisent les flèches permises.
\end{explain}
\end{document}
